\documentclass[10pt,a4paper]{amsart}
\usepackage[margin=19mm]{geometry}
\usepackage{amsmath,amssymb,mathtools}
\usepackage[scr=rsfs,cal=euler]{mathalfa}
\usepackage{xfrac}
\usepackage[unicode,hidelinks,bookmarks=false]{hyperref}
\newcommand{\ie}{i.e.\ }
\newcommand{\cf}{cf.\ }
\newcommand{\resp}{respectively\ }
\newcommand{\Subsection}{\subsection}
\providecommand{\nobreakdash}{\nobreak\hbox{-}}
\setlength{\emergencystretch}{2em}
\multlinegap0pt
\usepackage{amssymb}
\usepackage{xcolor}
\newtheorem{theorem}{Theorem}
\def\R{\mathcal{ R}}
\title{Generalization of Newton's minimal resistance problem to Riemannian surfaces}
\author{Rafael L\'opez}
\date{Source: arXiv:2605.27029v1 (CC BY 4.0)}
\begin{document}
\maketitle

\noindent\emph{Source and licence.} Generalization of Newton's minimal resistance problem to Riemannian surfaces. Version arXiv:2605.27029v1 (CC BY 4.0), distributed by arXiv under the Creative Commons Attribution 4.0 International licence, http://creativecommons.org/licenses/by/4.0/, as recorded in the arXiv licence metadata for this exact version and shown on the abstract page https://arxiv.org/abs/2605.27029v1. Full licence notice: \texttt{licenses/arxiv-2605.27029-CC-BY-4.0.txt}.\par\smallskip

\noindent\textit{Editorial scope.} Counted unit: the article's theorem t61 (source0.tex lines 665--674). The article's complete original proof of it is delivered with it (source0.tex lines 676--701, one \texttt{proof} environment of 26 lines). No mathematics is written, repaired or re-derived: the statement and the proof are the article's own lines, in the article's own notation, and no retained line is substituted or re-flowed.  Retained uncounted as the source-local context the proof needs: lines 304--306.  Nothing was omitted from any retained span, so the package carries no omission record.  No retained line contains a citation, so the wrapper carries no bibliography and no bibliography entry.  No retained line contains a figure environment, an image-inclusion command or drawing code, so no image is delivered and the package declares no asset dependency.  Editorial formatting only, in addition: an AMS \texttt{amsart} preamble that restates the article's own declarations and macros used by the retained lines, namely the environments \texttt{theorem} on separate counters, together with the standard packages those lines need.  The retained environments print in this document as Theorem 1 in order of appearance, whereas the article prints the same items with the numbering of its own full document.  The complete native source of the article as distributed in the arXiv e-print for this exact version is bundled unchanged as \texttt{sources/201464/source0.tex}.
\begin{equation}\label{r5}
\mathcal{R}=\frac{v_1-v_0}{1+k^2}.
\end{equation}

 \begin{theorem} \label{t61}
 Let  $A=(u_0, v_0)$ and $B=(u_1, v_1)$ be two points with $u_0 < u_1$ and $v_0<v_1$. If the angular amplitude $\Delta v = v_1 - v_0$ satisfies 
\begin{equation}
\Delta v \le \Phi(u_1) - \Phi(u_0),
\end{equation}
then the unique global minimizer of $\mathcal{R}$ in the class of admissible monotonic curves $\mathcal{C}_{0,1}^{m}$ is the loxodrome $\gamma_{lox}$ defined in geodesic coordinates by 
$$\Phi(u(v)) = k(v-v_0)+\Phi(u_0),$$
 where 
$$k = \frac{\Phi(u_1) - \Phi(u_0)}{\Delta v}.$$
\end{theorem}

\begin{proof}
We introduce the change of variables $y(v) = \Phi(u(v))$. Notice that the constraint $u' \ge 0$ is preserved as $y' \ge 0$. The resistance functional becomes 
\begin{equation}
\mathcal{R}[y] = \int_{v_0}^{v_1} \frac{1}{1 + y'^2}\, dv.
\end{equation}
The curve $\gamma_{lox}$ is, by definition, an extremal and it is immediate that $\gamma_{lox}$ connects $A$ to $B$ because $\Phi(u(v_0))=\Phi(u_0)$ and $\Phi(u(v_1))=k\Delta v+\Phi(u_0)=\Phi(u_1)$ by the definition of $k$. 

Let $g(p) = \frac{1}{1+p^2}$ be the integrand of $\mathcal{R}[y]$. To find the global minimum, we consider the lower convex envelope of $g(p)$ on the domain $p \in [0, \infty)$, denoted by $g^{**}(p)$. Since $g$ is concave on $[0,1)$ and convex on $(1,\infty)$, the function $g^{**}$ in the interval $[0,1]$ is the tangent line to $g$ from $(0, g(0)) = (0, 1)$ to $(1,g(1))=(1,\frac12)$; on $[1,\infty)$, we have $g^{**}=g$. This gives
\begin{equation}
g^{**}(p) = \begin{cases} 
1 - \frac{1}{2}p & \text{if } 0 \le p \le 1 \\
\frac{1}{1+p^2} & \text{if } p \geq 1 
\end{cases}
\end{equation}
Let $u\in \mathcal{C}_{0,1}^m$. Since $g\geq g^{**} $ and $g^{**}$ is convex, we can apply Jensen's inequality obtaining
\begin{equation}
 \begin{split}
\R[u]&= \int_{v_0}^{v_1} g(y')\, dv \ge \int_{v_0}^{v_1} g^{**}(y')\, dv \ge \Delta v\cdot g^{**} \left( \frac{1}{\Delta v} \int_{v_0}^{v_1} y'\, dv \right) \\
&= \Delta v\cdot g^{**} \left( \frac{1}{\Delta v} (\Phi(u_1)-\Phi(u_0)) \right)= \Delta v\cdot g^{**}(k)\\
&=\Delta v\cdot g(k),
\end{split}
\end{equation}
where in the last equality, we used the fact that $k\geq 1$ and the definition of $g^{**}$. Furthermore, the resistance of the loxodrome is given in \eqref{r5}, which yields $\R[\gamma_{lox}]=\frac{\Delta v}{1+k^2}=\Delta v\cdot g(k)$. Thus, $\R[u]\geq \R[\gamma_{lox}]$. 

To establish uniqueness, assume there is another curve $u=u(v)$ whose resistance is $\Delta v\cdot g(k)$. Then we have equality throughout the chain of inequalities. In particular, equality in Jensen's inequality holds if and only if the argument $y'(v)$ is constant almost everywhere. Thus $y'(v) = k$ identically and this proves that $u(v)$ is a loxodrome. By the boundary conditions at $v_0$ and $v_1$, $u(v)$ must be $\gamma_{lox}$. 
\end{proof}

\end{document}
